\labsection{2}{Arrays, multidimensional storage, and matrices}{sec:lab02}

\begin{sourcealignment}{Official Lab Manual --- Lab II}
\textbf{Objective.} Understand the array as a basic prerequisite for more complex data structures.

\textbf{Outcomes.} Explain how arrays are stored in memory and use arrays to solve programming problems.

\textbf{Problem directions.} Real-world scenarios involving arrays; multidimensional arrays; matrix-related problems such as matrix multiplication.
\end{sourcealignment}

\begin{labproject}{Process and multiply matrices}
\textbf{Coverage.} Chapter~\ref{ch:arrays}.\\
\textbf{Goal.} Multiply two matrices and explain what each of the three nested loops controls.

The lab example multiplies $A=\begin{psmallmatrix}1&2&3\\4&5&6\end{psmallmatrix}$ ($2\times3$) by $B=\begin{psmallmatrix}7&8\\9&10\\11&12\end{psmallmatrix}$ ($3\times2$). The multiplication refuses incompatible shapes before it starts:

\begin{lstlisting}
bool multiply(const int A[][MAXN], int aRows, int aCols,
              const int B[][MAXN], int bRows, int bCols,
              int P[][MAXN]) {
    if (aCols != bRows) return false;
    const int R = aRows, K = aCols, C = bCols;

    for (int i = 0; i < R; ++i)
        for (int j = 0; j < C; ++j) {
            P[i][j] = 0;
            for (int t = 0; t < K; ++t)
                P[i][j] += A[i][t] * B[t][j];
        }
    return true;
}
\end{lstlisting}

\begin{tryit}
Add two of the following: row sums, column sums, transpose, largest element, diagonal sum, or a check for whether a square matrix is symmetric.
\end{tryit}

\begin{workedexample}
For matrix multiplication, the result size is $R\times C$. Each output cell $P[i][j]$ accumulates products across the shared dimension $K$.
\end{workedexample}
\end{labproject}

\begin{verification}
\begin{itemize}
  \item Validate that the number of columns in $A$ equals the number of rows in $B$.
  \item Print both inputs and the output matrix with clear row/column formatting.
  \item Trace one output cell by hand and match it to the program result.
  \item Explain what each nested loop controls.
\end{itemize}
\end{verification}

